%! Elimination Matrices and Inverse Matrices — Unit 2, Lesson 2.2 — Group Activity
\documentclass[10pt]{article}
\usepackage{linalg-article}
\usepackage{linalg-boxes}
\newcommand{\TallMath}[1]{$\displaystyle #1\rule[-1.4em]{0pt}{3.2em}$}
\newcommand{\xx}{\mathbf{x}}
\newcommand{\bb}{\mathbf{b}}

\begin{document}

\pageheader{Unit 2, Lesson 2.2}{Group Activity}

\vspace{0.10in}

\begin{headlinebox}{blush}
\small\textbf{Build the Solving Machine.} A single elimination step is a matrix
$E=\begin{bsmallmatrix} 1&0\\-\ell&1 \end{bsmallmatrix}$; undoing it flips the sign
($E^{-1}=\begin{bsmallmatrix} 1&0\\+\ell&1 \end{bsmallmatrix}$). Chain the undo-steps far enough and
you get $A^{-1}$, the machine that solves $A\xx=\bb$ in one multiplication. Work with your partner
and \emph{show every step}.
\end{headlinebox}

\vspace{0.10in}

\begin{tcolorbox}[colback=white, colframe=black!40, title=\textbf{Tier R --- Remediate},
    fonttitle=\bfseries, arc=1.5mm, left=3mm, right=3mm, top=2mm, bottom=2mm, breakable]
Let $A=\begin{bmatrix} 2 & 5 \\ 4 & 12 \end{bmatrix}$. The pivot is $2$, so the multiplier is
$\ell=4\div2=2$.
\begin{enumerate}[leftmargin=1.7em, itemsep=8pt]
  \item Write the elimination matrix $E_{21}$ (put $-\ell$ in the lower-left of $I$).
        \; $E_{21}=\begin{bmatrix} 1 & 0 \\ \blank{0.9cm} & 1 \end{bmatrix}$ \vspace{0.2cm}
  \item Multiply $E_{21}A$ to reach upper-triangular $U$.
\begin{work}
  -2(2,5)+(4,12) &= (0,2) \\
  E_{21}A &= \begin{bsmallmatrix} 2&5\\0&2 \end{bsmallmatrix} = U
\end{work}
  \item Write $E_{21}^{-1}$, the matrix that \emph{undoes} the step. \; $E_{21}^{-1}=\begin{bmatrix} 1 & 0 \\ \blank{0.9cm} & 1 \end{bmatrix}$ \vspace{0.2cm}
  \item In one sentence: how are $E_{21}$ and $E_{21}^{-1}$ related? \writeline
\end{enumerate}
\end{tcolorbox}

\begin{tcolorbox}[colback=white, colframe=black!40, title=\textbf{Tier A --- Approaching Proficiency},
    fonttitle=\bfseries, arc=1.5mm, left=3mm, right=3mm, top=2mm, bottom=2mm, breakable]
\textbf{One machine, many orders.} A gift shop packs two box types. Each \textbf{Small} box holds
$1$ mug and $1$ candle; each \textbf{Large} box holds $2$ mugs and $3$ candles. With $s$ smalls and
$t$ larges, the totals are $A\begin{bsmallmatrix} s\\ t \end{bsmallmatrix}=\bb$ where
$A=\begin{bmatrix} 1 & 2 \\ 1 & 3 \end{bmatrix}$ (mug row, candle row).
\begin{enumerate}[leftmargin=1.7em, itemsep=8pt]
  \item Build the machine once: compute $ad-bc$ and write $A^{-1}$.
\begin{work}
  ad-bc &= 3-2 = 1 \\
  A^{-1} &= \begin{bsmallmatrix} 3&-2\\-1&1 \end{bsmallmatrix}
\end{work}
  \item \textbf{Monday's order:} $8$ mugs, $11$ candles. Use $\xx=A^{-1}\bb$ to find $s$ and $t$.
\begin{work}
  \begin{bsmallmatrix} 3&-2\\-1&1 \end{bsmallmatrix}\begin{bsmallmatrix} 8\\11 \end{bsmallmatrix}
      &= \begin{bsmallmatrix} 24-22\\-8+11 \end{bsmallmatrix}
       = \begin{bsmallmatrix} 2\\3 \end{bsmallmatrix}
\end{work}
  \item \textbf{Tuesday's order:} $10$ mugs, $14$ candles. Same machine --- find $s$ and $t$.
\begin{work}
  \begin{bsmallmatrix} 3&-2\\-1&1 \end{bsmallmatrix}\begin{bsmallmatrix} 10\\14 \end{bsmallmatrix}
      &= \begin{bsmallmatrix} 30-28\\-10+14 \end{bsmallmatrix}
       = \begin{bsmallmatrix} 2\\4 \end{bsmallmatrix}
\end{work}
  \item \emph{Interpret:} give each day's box counts, and say why building $A^{-1}$ once beats
        re-running elimination every day. \writeline
\end{enumerate}
\end{tcolorbox}

\begin{tcolorbox}[colback=white, colframe=black!40, title=\textbf{Tier E --- Extension},
    fonttitle=\bfseries, arc=1.5mm, left=3mm, right=3mm, top=2mm, bottom=2mm, breakable]
\textbf{(1) Gauss--Jordan.} Find $A^{-1}$ for $A=\begin{bmatrix} 1 & 2 \\ 3 & 7 \end{bmatrix}$ by
row-reducing $[\,A\mid I\,]$ down \emph{then} up until the left side is $I$.
\[
  \left[\begin{array}{cc|cc} 1 & 2 & 1 & 0 \\ 3 & 7 & 0 & 1 \end{array}\right]
\]
\begin{work}
  \text{row}_2-3\,\text{row}_1:\quad
      &\left[\begin{smallmatrix} 1&2&1&0\\0&1&-3&1 \end{smallmatrix}\right] \\
  \text{row}_1-2\,\text{row}_2:\quad
      &\left[\begin{smallmatrix} 1&0&7&-2\\0&1&-3&1 \end{smallmatrix}\right] \\
  A^{-1} &= \begin{bsmallmatrix} 7&-2\\-3&1 \end{bsmallmatrix}
\end{work}

\textbf{(2) No machine.} Try to invert $A=\begin{bmatrix} 2 & 3 \\ 4 & 6 \end{bmatrix}$. Compute
$ad-bc$, run one elimination step, and explain what goes wrong. Is $A$ \textbf{singular}? Connect it
to the zero-pivot breakdown from Lesson~2.1.
\begin{work}
  ad-bc &= 12-12 = 0 \\
  \text{row}_2-2\,\text{row}_1 &= (0,0)
\end{work}
\writelines{2}
\end{tcolorbox}

\end{document}
